% GENERATED FILE. Edit canonical lessons, then run npm run generate:books.
% canonical-source-hash: fnv1a64:d79e139ca574f80c
% canonical-lessons: FA-C109-abru, FA-C109-pust, FA-C109-ostokhan, FA-C109-farsh, FA-C109-mobl

\chapter{Eyebrow, Skin, Bone, Carpet, Sofa}
\label{ch:fa-eyebrow-skin-bone-carpet-sofa}
\hlchaptermodality{109}

\begin{chapteropening}
\textbf{By the end of this chapter:} \emph{I can say an eyebrow, the skin, a bone, a carpet and a sofa.}

Complete the last lesson of chapter 109: I can say an eyebrow, the skin, a bone, a carpet and a sofa.
\end{chapteropening}

\section[\texorpdfstring{\fa{ابرو} (abru)}{ابرو (abru)}]{\fa{ابرو} (abru) — an eyebrow}

\label{lesson:FA-C109-abru}



\begin{quote}
Before the new one: say the Persian for the chest, then the Persian for the waist, the back.
\end{quote}

\subsection*{You'll want to know: \fa{ابرو}}
\textbf{\fa{ابرو}} — \emph{abru} — \textquotedblleft{}an eyebrow\textquotedblright{}.

\begin{cousinweb}
\emph{abru} is a native Persian word, inherited rather than borrowed: it continues Middle Persian \emph{brūg}, \textquotedblleft{}eyebrow\textquotedblright{} (EtymDB). EtymDB takes it back to the Indo-European word behind English \emph{brow}.
\end{cousinweb}

\begin{grammarlens}[title={What you've built}]
One more word for this chapter.
\end{grammarlens}

\subsection*{Your turn}
\begin{itemize}
\raggedright
  \item \emph{Say it:} \emph{abru}
  \item \emph{Say it:} \emph{abru}, once more
  \item \emph{Say it:} say \emph{abru}
  \item \emph{Recall:} say the Persian for the chest, then the Persian for the waist, the back, then say \emph{abru} again
\end{itemize}

\begin{tcolorbox}[breakable,colback=teal!4,colframe=teal!35!black,title={Before you move on}]
What does \fa{ابرو} mean? (\textquotedblleft{}An eyebrow\textquotedblright{}.) Say it once more.
\end{tcolorbox}
\section[\texorpdfstring{\fa{پوست} (pust)}{پوست (pust)}]{\fa{پوست} (pust) — the skin}

\label{lesson:FA-C109-pust}



\begin{quote}
Before the new one: say the Persian for the waist, the back, then the Persian for an eyebrow.
\end{quote}

\subsection*{You'll want to know: \fa{پوست}}
\textbf{\fa{پوست}} — \emph{pust} — \textquotedblleft{}the skin\textquotedblright{}.

\begin{cousinweb}
\emph{pust} is a native Persian word, inherited rather than borrowed; it is usually traced to Middle Persian \emph{pōst}, \textquotedblleft{}skin, hide\textquotedblright{}.
\end{cousinweb}

\begin{grammarlens}[title={What you've built}]
One more word for this chapter.
\end{grammarlens}

\subsection*{Your turn}
\begin{itemize}
\raggedright
  \item \emph{Say it:} \emph{pust}
  \item \emph{Say it:} \emph{pust}, once more
  \item \emph{Say it:} say \emph{pust}
  \item \emph{Recall:} say the Persian for the waist, the back, then the Persian for an eyebrow, then say \emph{pust} again
\end{itemize}

\begin{tcolorbox}[breakable,colback=teal!4,colframe=teal!35!black,title={Before you move on}]
What does \fa{پوست} mean? (\textquotedblleft{}The skin\textquotedblright{}.) Say it once more.
\end{tcolorbox}
\section[\texorpdfstring{\fa{استخوان} (ostokhân)}{استخوان (ostokhân)}]{\fa{استخوان} (ostokhân) — a bone}

\label{lesson:FA-C109-ostokhan}



\begin{quote}
Before the new one: say the Persian for an eyebrow, then the Persian for the skin.
\end{quote}

\subsection*{You'll want to know: \fa{استخوان}}
\textbf{\fa{استخوان}} — \emph{ostokhân} — \textquotedblleft{}a bone\textquotedblright{}.

\begin{cousinweb}
\emph{ostokhân} is a native Persian word, inherited rather than borrowed; it is usually traced to Middle Persian \emph{astuxān}, \textquotedblleft{}bone\textquotedblright{}. It is usually set beside Greek \emph{osteon} and Latin \emph{os}, \textquotedblleft{}bone\textquotedblright{}, as one Indo-European word.
\end{cousinweb}

\begin{grammarlens}[title={What you've built}]
One more word for this chapter.
\end{grammarlens}

\subsection*{Your turn}
\begin{itemize}
\raggedright
  \item \emph{Say it:} \emph{ostokhân}
  \item \emph{Say it:} \emph{ostokhân}, once more
  \item \emph{Say it:} say \emph{ostokhân}
  \item \emph{Recall:} say the Persian for an eyebrow, then the Persian for the skin, then say \emph{ostokhân} again
\end{itemize}

\begin{tcolorbox}[breakable,colback=teal!4,colframe=teal!35!black,title={Before you move on}]
What does \fa{استخوان} mean? (\textquotedblleft{}A bone\textquotedblright{}.) Say it once more.
\end{tcolorbox}
\section[\texorpdfstring{\fa{فرش} (farsh)}{فرش (farsh)}]{\fa{فرش} (farsh) — a carpet}

\label{lesson:FA-C109-farsh}



\begin{quote}
Before the new one: say the Persian for the skin, then the Persian for a bone.
\end{quote}

\subsection*{You'll want to know: \fa{فرش}}
\textbf{\fa{فرش}} — \emph{farsh} — \textquotedblleft{}a carpet\textquotedblright{}.

\begin{cousinweb}
\emph{farsh} is borrowed from Arabic \emph{farš}, \textquotedblleft{}something spread out, carpet\textquotedblright{}. Its Arabic root is \emph{f-r-š} (Lane).
\end{cousinweb}

\begin{grammarlens}[title={What you've built}]
One more word for this chapter.
\end{grammarlens}

\subsection*{Your turn}
\begin{itemize}
\raggedright
  \item \emph{Say it:} \emph{farsh}
  \item \emph{Say it:} \emph{farsh}, once more
  \item \emph{Say it:} say \emph{farsh}
  \item \emph{Recall:} say the Persian for the skin, then the Persian for a bone, then say \emph{farsh} again
\end{itemize}

\begin{tcolorbox}[breakable,colback=teal!4,colframe=teal!35!black,title={Before you move on}]
What does \fa{فرش} mean? (\textquotedblleft{}A carpet\textquotedblright{}.) Say it once more.
\end{tcolorbox}
\section[\texorpdfstring{\fa{مبل} (mobl)}{مبل (mobl)}]{\fa{مبل} (mobl) — a sofa}

\label{lesson:FA-C109-mobl}



\begin{quote}
Before the new one: say the Persian for a bone, then the Persian for a carpet.
\end{quote}

\subsection*{You'll want to know: \fa{مبل}}
\textbf{\fa{مبل}} — \emph{mobl} — \textquotedblleft{}a sofa\textquotedblright{}.

\begin{cousinweb}
\emph{mobl} is a loan, usually traced to French \emph{meuble}, \textquotedblleft{}piece of furniture\textquotedblright{}.
\end{cousinweb}

\begin{grammarlens}[title={What you've built}]
One more word for this chapter.
\end{grammarlens}

\subsection*{Your turn}
\begin{itemize}
\raggedright
  \item \emph{Say it:} \emph{mobl}
  \item \emph{Say it:} \emph{mobl}, once more
  \item \emph{Say it:} say \emph{mobl}
  \item \emph{Recall:} say the Persian for a bone, then the Persian for a carpet, then say \emph{mobl} again
\end{itemize}

\begin{tcolorbox}[breakable,colback=teal!4,colframe=teal!35!black,title={Before you move on}]
What does \fa{مبل} mean? (\textquotedblleft{}A sofa\textquotedblright{}.) Say it once more.
\end{tcolorbox}
